% Options for packages loaded elsewhere
\PassOptionsToPackage{unicode}{hyperref}
\PassOptionsToPackage{hyphens}{url}
\documentclass[
  ignorenonframetext,
]{beamer}
\newif\ifbibliography
\usepackage{pgfpages}
\setbeamertemplate{caption}[numbered]
\setbeamertemplate{caption label separator}{: }
\setbeamercolor{caption name}{fg=normal text.fg}
\beamertemplatenavigationsymbolsempty
% remove section numbering
\setbeamertemplate{part page}{
  \centering
  \begin{beamercolorbox}[sep=16pt,center]{part title}
    \usebeamerfont{part title}\insertpart\par
  \end{beamercolorbox}
}
\setbeamertemplate{section page}{
  \centering
  \begin{beamercolorbox}[sep=12pt,center]{section title}
    \usebeamerfont{section title}\insertsection\par
  \end{beamercolorbox}
}
\setbeamertemplate{subsection page}{
  \centering
  \begin{beamercolorbox}[sep=8pt,center]{subsection title}
    \usebeamerfont{subsection title}\insertsubsection\par
  \end{beamercolorbox}
}
% Prevent slide breaks in the middle of a paragraph
\widowpenalties 1 10000
\raggedbottom
\AtBeginPart{
  \frame{\partpage}
}
\AtBeginSection{
  \ifbibliography
  \else
    \frame{\sectionpage}
  \fi
}
\AtBeginSubsection{
  \frame{\subsectionpage}
}
\usepackage{iftex}
\ifPDFTeX
  \usepackage[T1]{fontenc}
  \usepackage[utf8]{inputenc}
  \usepackage{textcomp} % provide euro and other symbols
\else % if luatex or xetex
  \usepackage{unicode-math} % this also loads fontspec
  \defaultfontfeatures{Scale=MatchLowercase}
  \defaultfontfeatures[\rmfamily]{Ligatures=TeX,Scale=1}
\fi
\usepackage{lmodern}
\ifPDFTeX\else
  % xetex/luatex font selection
\fi
% Use upquote if available, for straight quotes in verbatim environments
\IfFileExists{upquote.sty}{\usepackage{upquote}}{}
\IfFileExists{microtype.sty}{% use microtype if available
  \usepackage[]{microtype}
  \UseMicrotypeSet[protrusion]{basicmath} % disable protrusion for tt fonts
}{}
\makeatletter
\@ifundefined{KOMAClassName}{% if non-KOMA class
  \IfFileExists{parskip.sty}{%
    \usepackage{parskip}
  }{% else
    \setlength{\parindent}{0pt}
    \setlength{\parskip}{6pt plus 2pt minus 1pt}}
}{% if KOMA class
  \KOMAoptions{parskip=half}}
\makeatother
\setlength{\emergencystretch}{3em} % prevent overfull lines
\providecommand{\tightlist}{%
  \setlength{\itemsep}{0pt}\setlength{\parskip}{0pt}}
\usepackage{bookmark}
\IfFileExists{xurl.sty}{\usepackage{xurl}}{} % add URL line breaks if available
\urlstyle{same}
\hypersetup{
  hidelinks,
  pdfcreator={LaTeX via pandoc}}

\author{\texorpdfstring{}{}}
\date{}

\begin{document}

\begin{frame}{The Bullionist Controversy and the Semiotics of Value}
\protect\phantomsection\label{sec-bullionist-controversy}
The aftermath of the 1797 Bank Restriction Act birthed the Bullionist
Controversy, a political and economic debate that raged throughout the
early 19th century and culminated in the 1810 Bullion Report presented
to Parliament. Crucially, this economic debate mapped closely onto the
semiotic and epistemological battles William Blake was waging in his
prophetic poetry. The central issue of the Controversy was whether the
rising price of bullion resulted from the over-issuance of paper money
(the ``Bullionist'' position), or if the paper money was maintaining its
value while gold itself was appreciating due to wartime scarcity (the
``Anti-Bullionist'' position) {[}@fetter1965development{]}.

\textbf{Bullionists (Ricardo and the metallic anchor).} The Bullionists,
led by figures like David Ricardo and embodied in the 1810 Bullion
Report, argued that value must refer back to a physical, metallic
anchor. Ricardo's \emph{The High Price of Bullion, a Proof of the
Depreciation of Bank Notes} (1810) provided the intellectual foundation:
the premium on gold bullion over the nominal value of banknotes
constituted empirical proof that the Bank had depreciated the currency
through overissuance {[}@ricardo1810high{]}. Henry Thornton, in his
\emph{An Enquiry into the Nature and Effects of the Paper Credit of
Great Britain} (1802), offered a more nuanced but equally comprehensive
analysis, demonstrating that the Bank's credit expansion operated with
radical independence from any metallic constraint, introducing systemic
instability into all commodity prices {[}@thornton1802nature{]}. For the
Bullionists, the post-1816 pound was effectively anchored to the
sovereign's 113 grains of fine gold, and any deviation from that anchor
was considered a species of national fraud.

\textbf{Anti-Bullionists (the Bank's directors).} The Anti-Bullionists,
representing the Bank of England's directors and their parliamentary
allies, paradoxically argued for a more abstract, relational conception
of value. They claimed that their paper notes represented not a fixed
weight of metal, but the ``needs of trade''---the aggregate commercial
demand of the nation. This was an early doctrine of paper credit and
convertibility: that the Bank's notes could circulate so long as
confidence and state backing sustained them
{[}@fetter1965development{]}. From the perspective of Blakean cosmology,
this Anti-Bullionist position occupied a sophisticated space: it
acknowledged the constructed nature of monetary value, but deployed this
acknowledgment to entrench the oligarchic power of the Bank's governors
rather than to liberate humanity from quantification.

\begin{block}{Blake's Refusal: Newton's Mint and Los's Forge}
\protect\phantomsection\label{blakes-refusal-newtons-mint-and-loss-forge}
In this debate, Blake's position is complex and irreducible to either
camp; he frames the crisis not as macroeconomic policy, but as a
political theology pitting \textbf{Newton's Mint} against \textbf{Los's
Forge}. He was certainly not a Bullionist---he abhorred the ``chimeric
regime'' of the gold standard and the worship of the ``guinea'' as a
false idol, declaring in his marginalia to Reynolds that to a miser, ``a
guinea is more beautiful than the sun'' {[}@blake1808annotations{]}.
Yet, he was equally critical of the Anti-Bullionist reality: a financial
oligarchy manipulating abstract paper contracts (``allegoric riches'')
to extract real labor from the poor. \emph{Jerusalem} Plate 27 targets
that dual exploitation at length (see \hyperlink{sec-allegoric-riches}{§
The fixing of labour and allegoric riches}) {[}@erdman1988complete;
@sklar2011blake{]}.

Both the Bullionists and the Anti-Bullionists were trapped in
\emph{Urizenic} logic. One worshipped the atomic, finite particle (the
gold coin); the other worshipped the algorithmic void of credit
(fractional reserve banking). Against both, Blake posited Los's Forge:
the imaginative, artisanal production of eternal forms through immediate
physical labor. The material toll of this macroeconomic abstraction on
Los's Forge was substantive. As the Restriction era dragged on, the
economic pressures of currency instability and inflation severed the
working classes from the means of artistic production. By 1818,
overwhelmed by this fiat malaise, Blake virtually ceased printing his
illuminated books---producing only 14 copies of four titles
post-1795---and pivoted to commercial graphic art to survive
{[}@makdisi2003william{]}. The abstract ledgers of Threadneedle Street
had literally cooled the fires of the prophetic forge.

Blake's \emph{Fourfold Vision} (detailed under
\hyperlink{sec-urizenic-materialism}{§ Urizenic materialism}) entirely
circumvents this dichotomy. For Blake, true value is not found in the
metallic anchor nor in the abstract paper calculus, but in the
generative, qualitative labor of the human imagination. The Bullionist
Controversy represented the ultimate crisis of capitalist semiotics
{[}@fetter1965development{]}---a society arguing over whether the
\emph{sign} (paper) should be bound to the \emph{referent} (gold), while
entirely ignoring the \emph{creator} of value (the human
artist/laborer). As Mary Poovey demonstrates in her analysis of ``genres
of the credit economy,'' the entire semiotic infrastructure of
eighteenth-century finance was built upon an irresolvable tension
between the promise of face value and the reality of market value---a
tension that Blake's prophetic works expose as symptomatic of the deeper
spiritual crisis of quantification itself {[}@poovey2008genres{]}. The
Controversy was eventually decided not by intellectual resolution but by
legislative force. Parliament had already codified gold-centred coinage
and token silver during Restriction in the \textbf{1816} Coinage Act
{[}@redish1990evolution{]}; \textbf{Peel's Act} of \textbf{1819} then
committed Britain to \textbf{resumption} of Bank-note convertibility
into specie by \textbf{1821}, closing the twenty-four-year arc of the
Restriction period {[}@fetter1965development{]} and binding the restored
metallic order to that statutory frame.
\end{block}
\end{frame}

\end{document}
